\talkslide{rollouts}{Putting it all together}{Timing: provisional}{Follow one seed through the computation.}{
\jointtrainingbuild{0}
}
\speak{Follow one seed through the computation. Two proposed prompts create two branches. For each faithful prompt, generate two defender responses and rank them. Those within-branch pairs train the defender. Then compare the preferred response from each branch under the attacker's criterion to obtain an attacker pair. If only one attack is faithful, that rule decides the attacker winner directly. If neither is faithful, skip the seed. The diagram shows the case where both attacks are faithful. The two response pairs feed the defender loss shown at the lower left, averaging across attacks and seeds. The attack pair feeds the attacker loss at the lower right. Both use the same reference-relative score-gap form as the preceding DPO slides; the fixed conditioning variables are suppressed in this compact view. The defender judge is drawn as one grouped stage containing two independent comparisons. We refresh the generators and repeat. The practical recipe samples from exponential moving averages of both models. The paper writes its losses relative to reference policies. The released Qwen and Llama DPO configurations also use EMA log probabilities inside the loss. This is an implementation choice, not a fixed-reference claim.}
